\section{Introduction}\label{sec:intro}

Discrete phase models support pinned defects, saddle-mediated transitions, and competing escape paths. Finding one path with a low energy maximum gives an upper bound on an activation barrier. Proving that this path is globally optimal requires excluding every less expensive path, including those that leave a proposed reaction surface or create additional defects. On an infinite lattice, even a very accurate finite-box saddle leaves an exterior verification problem. A small residual at the interior vertices does not control the residual created at the interface of an extension.

This paper addresses that distinction for a neutral dipole in the square-lattice XY model with the exact interaction $1-\cos t$. We work in the homogeneous energy Hilbert space $\X$ defined in \cref{sec:model}. Our central result identifies the annihilation barrier over \emph{all} strongly $\X$-continuous paths to any physical uniform lift. It is a variational assertion about the full infinite lattice, rather than a minimum within a frozen two-coordinate family.

\subsection{Results and their scope}

For a dipole with charges $+1$ and $-1$ at plaquettes $(0,0)$ and $(12,0)$, the computation certifies a minimum $m_0$ and two reflected asymmetric index-one saddles $s_+$ and $s_-$. Their common energy is denoted by $h$. An open bounded set $\mathcal D\subset\X$ contains $m_0$, excludes every annihilation endpoint from its closure, and satisfies $\E\ge h$ on its entire boundary. A separately validated path from $m_0$ to an endpoint has maximum $h$. Consequently,
\begin{equation}\label{eq:intro-barrier}
 B_0(m_0)=h-\E(m_0),\qquad
 0.00567867<B_0(m_0)<0.00570116.
\end{equation}
Here the coupling is one, and the numbers are dimensionless. The minimum and saddle assertions refer to general perturbations in $\X$; no symmetry restriction or positivity of every core bond is imposed. The meaning of ``all paths'' includes the specified strong topology and the endpoint set in \cref{eq:barrier}.

The geometric argument becomes reusable through \cref{thm:certificate}. Its hypotheses are a finite collection of rigorous inequalities, together with analytic formulas controlling the exterior. The conditions cover full dual residuals, signed-bond coercivity, saddle index, exact cutting planes, transverse fibers, all chart faces and seams, and endpoint attachment of the upper path. A certificate that fails one test is inconclusive; it does not prove instability or absence of a lower route. The verification format permits adaptation to another instance without asserting that the search for a successful certificate always terminates.

Three analytic extensions complement this fixed example. First, \cref{thm:green} gives an explicit Fourier estimate for the orthogonal projection onto lattice gradients. Combined with the stationary equation, it yields an inverse-radius energy tail. Its constants also control the already validated upper path. Second, \cref{thm:local-class,thm:class-termination} reuse the fixed
communication height for an explicit infinite class of initial fields.
The class is generated by compact rational blocks with a geometric
norm budget and a prescribed spatial support sequence. It allows
noncompact perturbations without imposing symmetry. A subcritical
connection to $m_0$ transfers the full path-space certificate, and a
constructive energy enclosure terminates for every rational accuracy.
This is a consequence of the existing global certificate, rather than a
search for all mechanisms. The general refinement procedure is proved;
the new code checks its constants and representative finite prefixes.
Third, \cref{thm:robustness} treats the infinite-dimensional family of interactions
\begin{equation}\label{eq:intro-perturbation}
 E_b(\theta)=\sum_e(1+b_e)(1-\cos D_e\theta),\qquad
 b\in\ell^1(\mathcal E),\quad\nrm b_1\le 1/500.
\end{equation}
The minimum and the two distinct index-one saddles continue, the minimum retains its charges, and
\begin{equation}\label{eq:intro-robust}
 |B_b(m_b)-B_0(m_0)|\le2\nrm b_1.
\end{equation}
Thus the perturbed global barrier remains strictly positive. The same coupling class admits a uniform inverse-radius tail for the three continued critical states. We also prove local exponential stability for the specified energy-space gradient flow. We do not claim that a continued saddle necessarily realizes the perturbed minimax, or that a mechanism switch has been established.

\subsection{Relation to prior work}

Certification of critical points of cosine energy landscapes is established. Mehta, Hauenstein, and Wales \cite{Mehta2014} used algebraic certification for finite XY systems with periodic or antiperiodic boundaries. More recently, Ayala, Choksi, and Wirth \cite{Ayala2026} validated finite atomistic configurations and transition-state candidates. These results show why neither an XY critical point certificate nor a validated finite-dimensional saddle should be presented as the new conclusion here.

Metastability and depinning for square-lattice systems with periodic interactions, including XY models, were studied variationally by Alicandro, De Luca, Garroni, and Ponsiglione \cite{Alicandro2014}. They established metastable configurations and derived vortex dynamics through minimizing movements and continuum-limit arguments. Here the lattice spacing and the vortex separation are fixed, and the quantity to be certified is the global communication height from a specified metastable state over all strong energy-space escape paths.

For lattice defects, Hudson and Ortner \cite{Hudson2015} obtained strongly stable multi-dislocation configurations from an isolated-core stability hypothesis and sufficient separation. We verify a particular dipole equilibrium and a global escape barrier at a fixed separation. Hudson \cite{Hudson2017} derived asymptotic energy barriers in a finite-domain, piecewise-quadratic dislocation model with a specified single-hop path class. Garroni, Petrache, and Spadaro \cite{Garroni2025} proved clearing-out results for global minimizers of discrete energies with topological singularities. Global minimization and pinned local metastability ask different questions; a clearing-out theorem does not supply a continuous path with a prescribed ceiling from a given metastable state.

Validated transition paths also have direct precedents. Breden and Kuehn \cite{Breden2019} rigorously validated saddle connections in a stochastic transition-path setting. Liu, Chen, and Ortner \cite{Liu2024} developed stability theory for minimum energy paths in Hilbert spaces under explicit nondegeneracy hypotheses. Berglund, Medvedev, and Simpson \cite{Berglund2025} analyzed metastability and relevant transition barriers in the finite one-dimensional nearest-neighbor Kuramoto model, also a cosine system. These works preclude a general claim of priority for a saddle-height formula, for infinite-dimensional path theory, or for rigorous transition geometry. Our lower bound instead comes from a complete separating boundary in the energy space of a two-dimensional infinite lattice. It does not assume that a previously selected path is a minimum energy path.

Several ingredients of the certificate also have established analogues. A posteriori existence arguments based on computable residuals and stability bounds have a substantial history; see Ortner \cite{Ortner2009} for an inverse-function framework for nonmonotone problems. Signed Laplacian positivity and inertia can be expressed through Schur complements and effective resistances \cite{Chen2020}. Cubical sublevel computations and their verification are treated by D\l{}otko and Wanner \cite{Dlotko2018}.

Far-field estimates and boundary approximations were developed by Ehrlacher, Ortner, and Shapeev \cite{Ehrlacher2016}. Braun, Hudson, and Ortner \cite{Braun2022} obtained systematic far-field expansions using continuum correctors and discrete multipoles; higher-order boundary conditions were subsequently developed by Braun, Ortner, Wang, and Zhang \cite{Braun2025} and by Wei, Braun, Wang, and Zhang \cite{Wei2026}. For the present certificate we require explicit constants in matching bond conventions for the critical states and for an entire upper path.

The contribution is therefore a specific conjunction: the exact cosine interaction on the infinite square lattice; unrestricted strong energy-space escape paths; a matched global lower and upper barrier bound; checkable hypotheses exposing the infinite-dimensional reduction; a terminating reuse procedure on a specified infinite input class; and a quantitative perturbation theorem with effective tails in the same coupling class. We do not infer from this comparison that earlier methods are incapable of adaptation, nor assert universal bibliographic priority.

\subsection{Verification record and organization}

The numerical theorem uses frozen proposals and rigorous verification, with the three source archives indexed in the supplementary record. A fresh execution of the standard supervisor reproduced every phase it invokes, without using cached validated outputs as mathematical premises. This is a reproduction of the same verification software, supported by analytic review; it is not an independent interval implementation or a proof-assistant formalization. The reproducibility record distinguishes these levels of evidence.

\Cref{sec:model} fixes the topology and normalization. \Cref{sec:certificate} gives the reusable a posteriori theorem, and \cref{sec:instance} records its verified instance. \Cref{sec:green,sec:local-class,sec:robustness} prove the effective-tail, terminating local-class, and coupling-robustness extensions. \Cref{sec:discussion} discusses algorithmic scope and the next certification problem.
